%==============================================================================
%  The Distance Route: split recovery from a sub-sampled double-centred
%  distance matrix.
%
%  SELF-CONTAINED.  Own preamble, own bibliography (refs.bib, in this
%  directory).  Build from inside distance-route/ with
%
%      latexmk -pdf -interaction=nonstopmode distance_route.tex
%
%  It shares no \input, no label namespace and no .bib file with thesis_v9.tex,
%  and it must never be added to the master's \input list without first
%  renaming its labels -- every label here is prefixed d: for that reason.
%  Every macro is \providecommand'd so that an \input into a document which
%  already defines \Dmat/\Bop/... is a no-op rather than a clash.
%
%  [Reviewer Note: this document deliberately does NOT reuse the machinery of
%   the deleted sections/appendix-G.tex.  That route factored B through the
%   Bandelt--Dress canonical decomposition, stated its only closed-form gap for
%   balanced splits, and carried a per-node cancellation hypothesis that the
%   route's own referee report records as false and removed.  The structure
%   theorems below (Theorem 1 and Theorem 2) are proved here from the
%   definition of an additive tree metric and from the Laplacian of the tree.]
%==============================================================================
\documentclass[11pt]{article}
\usepackage[margin=1in]{geometry}
\usepackage{amsmath,amssymb,amsthm,mathtools}
\usepackage{enumitem}
\usepackage{booktabs}
\usepackage{hyperref}
\usepackage{cleveref}

\crefname{assumption}{Assumption}{Assumptions}
\Crefname{assumption}{Assumption}{Assumptions}

\bibliographystyle{plain}

%------------------------------ theorem environments --------------------------
\theoremstyle{plain}
\newtheorem{theorem}{Theorem}
\newtheorem{lemma}{Lemma}
\newtheorem{proposition}{Proposition}
\newtheorem{corollary}{Corollary}

\theoremstyle{definition}
\newtheorem{definition}{Definition}
\newtheorem{assumption}{Assumption}
\renewcommand{\theassumption}{D\arabic{assumption}}

\theoremstyle{remark}
\newtheorem{remark}{Remark}

%------------------------------ notation --------------------------------------
\providecommand{\Dmat}{\mathcal{D}}          % leaf-to-leaf distance matrix
\providecommand{\Bop}{\mathcal{B}}           % centred distance operator H D H
\providecommand{\Mop}{\mathcal{M}}           % -Bop, positive semidefinite
\providecommand{\uone}{u^{(1)}}
\providecommand{\huone}{\hat u^{(1)}}
\providecommand{\1}{\mathbf{1}}
\providecommand{\E}{\mathbb{E}}
\providecommand{\Prob}{\mathbb{P}}
\providecommand{\R}{\mathbb{R}}
\providecommand{\norm}[1]{\left\lVert #1 \right\rVert}
\providecommand{\opnorm}[1]{\norm{#1}_{2}}
\providecommand{\hp}{with high probability}
\newcommand{\dmax}{d_{\max}}
\newcommand{\est}{e_{*}}
\newcommand{\wst}{w_{*}}
\newcommand{\Ast}{A_{*}}
\newcommand{\Bst}{B_{*}}
\newcommand{\ehat}{\hat{e}}

\title{The Distance Route:\\
       split recovery from a sub-sampled double-centred distance matrix}
\author{}
\date{}

\begin{document}
\maketitle

\begin{abstract}
Let $\Dmat$ be the matrix of path distances between the $m$ leaves of a weighted
tree and let $\Bop=H\Dmat H$ be its double centring, $H=I-\tfrac1m\1\1^\top$. We
give a self-contained account of why the sign pattern of the leading-magnitude
eigenvector of $\Bop$ is a bipartition of the leaves cut by a single edge of the
tree, and of how much of the matrix may be discarded before that sign pattern
changes. Two exact structure theorems carry the argument. The first writes
$\Bop$ as a negatively weighted sum of rank-one split projectors,
$\Bop=-\sum_{e}w_e\,\ehat_e\ehat_e^\top$ with $w_e=2\ell_e a_e b_e/m$; the second
identifies $-\tfrac12\Bop$ with the Moore--Penrose pseudoinverse of the tree
Laplacian Schur-complemented onto the leaves. The first supplies the spectral
gap, and an explicit preference for balanced bipartitions that the Laplacian
operator does not share; the second supplies validity, with no quantitative
hypothesis at all. Against these we run a sampling model in which each pair is
observed independently with probability $p$: an operator-norm bound for the
inverse-probability estimator, a Davis--Kahan step, and a one-step refinement
whose per-leaf statistic concentrates by a scalar Bernstein inequality. The
conclusion is that $p=\Theta(\log m/m)$ suffices to reproduce exactly the
partition the fully observed operator returns.
\end{abstract}

%==============================================================================
\section{Setting}\label{d:sec:setup}
%==============================================================================

Let $T$ be a tree with leaf set $X$, $|X|=m\ge 4$, internal vertex set $R$, edge
set $E$, and a strictly positive length $\ell_e$ on each edge. For $i,j\in X$
let $\Dmat_{ij}$ be the total length of the unique $i$--$j$ path, and write
\begin{equation}\label{d:eq:dmax}
  \dmax:=\max_{i,j\in X}\Dmat_{ij}.
\end{equation}
Let $H:=I-\tfrac1m\1\1^\top$ be the centring projector and set
\begin{equation}\label{d:eq:bdef}
  \Bop:=H\,\Dmat\,H,\qquad \Mop:=-\Bop .
\end{equation}
The object of study is the \emph{sign rule}: take $\uone$, a unit eigenvector of
$\Bop$ belonging to the eigenvalue of largest absolute value, and partition $X$
by $\operatorname{sign}(\uone)$.

\begin{remark}[Sign convention]\label{d:rem:sign}
\Cref{d:cor:nsd} below shows $\Bop\preceq0$, so the eigenvalue of largest
absolute value is the \emph{most negative} one, and $\uone$ is the top
eigenvector of $\Mop\succeq0$. The literature that uses this operator describes
$\uone$ as ``leading'' \cite{aizenbud2023spectral}; on $\Bop$ that word must be
read as largest in magnitude, not largest algebraically, since the largest
algebraic eigenvalue of $\Bop$ is the eigenvalue $0$ carried by $\1$.
\end{remark}

Throughout, $A_e\subset X$ and $B_e=X\setminus A_e$ denote the two leaf sets
obtained by deleting edge $e$ from $T$, with $a_e=|A_e|$ and $b_e=|B_e|$; both
are non-empty. A bipartition of $X$ of this form is called a \emph{split of
$T$}. For $A\subseteq X$ we write $\1_A$ for its indicator vector.

%==============================================================================
\section{The split-covariance identity}\label{d:sec:structure}
%==============================================================================

The first structure theorem needs nothing beyond the definition of an additive
metric and one line of algebra.

\begin{lemma}[Path decomposition]\label{d:lem:path}
For $e\in E$ put $S_e:=\1_{A_e}\1_{B_e}^\top+\1_{B_e}\1_{A_e}^\top$. Then
\begin{equation}\label{d:eq:path}
  \Dmat=\sum_{e\in E}\ell_e\,S_e .
\end{equation}
\end{lemma}

\begin{proof}
$T$ is a tree, so the $i$--$j$ path is unique and
$\Dmat_{ij}=\sum_{e\in E}\ell_e\,\mathbf{1}\{e\text{ lies on the }i\text{--}j
\text{ path}\}$. Deleting $e$ splits $T$ into the two components containing
$A_e$ and $B_e$ respectively, and $e$ lies on the $i$--$j$ path exactly when $i$
and $j$ fall in different components, i.e.\ exactly when $(S_e)_{ij}=1$.
\end{proof}

\begin{lemma}[Rank-one centring]\label{d:lem:centring}
Let $A\subseteq X$ with $a=|A|\in\{1,\dots,m-1\}$, $b=m-a$, and set
$\tilde u:=H\1_A$. Then
\begin{equation}\label{d:eq:centring}
  H\bigl(\1_A\1_{B}^\top+\1_{B}\1_A^\top\bigr)H=-2\,\tilde u\tilde u^\top,
  \qquad
  \norm{\tilde u}^2=\frac{ab}{m},
\end{equation}
where $B=X\setminus A$. Moreover $\tilde u_i=b/m$ for $i\in A$ and
$\tilde u_i=-a/m$ for $i\in B$, so with $\ehat:=\tilde u/\norm{\tilde u}$,
\begin{equation}\label{d:eq:deloc}
  \ehat_i=\sqrt{\tfrac{b}{am}}\ \ (i\in A),\qquad
  \ehat_i=-\sqrt{\tfrac{a}{bm}}\ \ (i\in B),\qquad
  \min_{i}|\ehat_i|=\sqrt{\frac{\min(a,b)}{m\,\max(a,b)}} .
\end{equation}
\end{lemma}

\begin{proof}
$H\1=0$ and $\1_B=\1-\1_A$, hence $H\1_B=-\tilde u$. Therefore
$H(\1_A\1_B^\top+\1_B\1_A^\top)H
 =(H\1_A)(H\1_B)^\top+(H\1_B)(H\1_A)^\top
 =-2\tilde u\tilde u^\top$.
The entries of $\tilde u=\1_A-\tfrac{a}{m}\1$ are as stated, so
$\norm{\tilde u}^2=a(b/m)^2+b(a/m)^2=ab(a+b)/m^2=ab/m$, and dividing gives
\eqref{d:eq:deloc}.
\end{proof}

Write $\tilde u_e:=H\1_{A_e}$ and $\ehat_e:=\tilde u_e/\norm{\tilde u_e}$.
Note that $\operatorname{sign}(\ehat_e)$ \emph{is} the split $(A_e,B_e)$, with no
zero coordinate.

\begin{theorem}[Split covariance]\label{d:thm:splitcov}
\begin{equation}\label{d:eq:splitcov}
  \Bop=-\sum_{e\in E}w_e\,\ehat_e\ehat_e^\top,
  \qquad
  w_e:=\frac{2\,\ell_e\,a_e\,b_e}{m}\;>\;0 .
\end{equation}
\end{theorem}

\begin{proof}
Apply $H(\cdot)H$ to \eqref{d:eq:path} and use \eqref{d:eq:centring} term by
term: $H(\ell_eS_e)H=-2\ell_e\tilde u_e\tilde u_e^\top
=-2\ell_e\norm{\tilde u_e}^2\ehat_e\ehat_e^\top=-w_e\ehat_e\ehat_e^\top$.
\end{proof}

\Cref{d:thm:splitcov} is the exact and finite counterpart, for a tree metric, of
a canonical split decomposition \cite{bandeltdress1992canonical}: the split
system of $T$ is its edge set, so no limiting or canonical-form argument is
required.

\begin{corollary}[Definiteness, rank, trace]\label{d:cor:nsd}
$\Bop\preceq0$ with $\Bop\1=0$; $\operatorname{rank}\Bop=m-1$, so the inertia of
$\Bop$ is $(0,1,m-1)$; and
\begin{equation}\label{d:eq:trace}
  \sum_{e\in E}w_e=\operatorname{tr}\Mop=\frac{1}{m}\,\1^\top\Dmat\,\1 .
\end{equation}
\end{corollary}

\begin{proof}
Negative semidefiniteness and $\Bop\1=0$ are immediate from
\eqref{d:eq:splitcov} and $\ehat_e^\top\1=0$. For the rank: the $m$ pendant
edges contribute the vectors $\tilde u_{\{i\}}=He_i$, $i\in X$, which span
$\1^{\perp}$, a space of dimension $m-1$; since every $w_e>0$ the quadratic form
$x^\top\Mop x=\sum_ew_e\langle \ehat_e,x\rangle^2$ vanishes on $\1^{\perp}$ only
at $0$. For the trace, each $\ehat_e$ is a unit vector, so
$\operatorname{tr}\Mop=\sum_ew_e$; independently
$\operatorname{tr}(H\Dmat H)=\operatorname{tr}(\Dmat H)
 =\operatorname{tr}\Dmat-\tfrac1m\1^\top\Dmat\1=-\tfrac1m\1^\top\Dmat\1$
because $\Dmat$ has zero diagonal.
\end{proof}

\begin{remark}[Inertia of $\Dmat$ itself]\label{d:rem:inertia}
\Cref{d:cor:nsd} is consistent with, and is the centred form of, the classical
statement that the distance matrix of a tree on $m$ vertices has exactly one
positive and $m-1$ negative eigenvalues. For \emph{weighted} trees --- the case
here --- the citable statement is the $s=1$ specialisation of Bapat's
matrix-weight determinant theorem \cite{bapat2006determinant}; Graham and
Pollak proved the determinant
$\det\Dmat=(-1)^{m-1}(m-1)2^{m-2}$ for the unweighted case
\cite{grahampollak1971addressing}, and Graham and Lov\'asz the associated
polynomials \cite{grahamlovasz1978polynomials}. The wider family of exact
relations between the distance matrix of a tree and its Laplacian, including the
closed form of $\Dmat^{-1}$, is collected in
\cite{bapat2005distance,bapat2014graphs}.
% [Reviewer Note: the exponent is m-2.  Several secondary sources print
%  2^{m-1}; that is a transcription error.  A second trap in the same
%  literature: the vector tau = 2*1 - d appearing in the inverse formula of
%  bapat2005distance uses the UNWEIGHTED degree, not the weighted degree.]
\end{remark}

\begin{remark}[$\Bop$ is the classical scaling operator]\label{d:rem:mds}
$G=-\tfrac12\Bop$ is the Gram matrix of classical multidimensional scaling
applied to $\Dmat$ \cite{young1938discussion,torgerson1952mds,gower1966distance},
so $\operatorname{sign}(\uone)$ is the sign of the first principal coordinate of
the leaves. \Cref{d:cor:nsd} states that this Gram matrix is positive
semidefinite for every weighted tree, which is the assertion that the tree metric
$\Dmat$ --- not its square --- is a squared Euclidean distance, i.e.\ that
$\sqrt{\Dmat}$ embeds isometrically in $\R^{m-1}$
\cite{schoenberg1935remarks,gower1985properties}; that phylogenetic distance
matrices have this property, and the elementary reason for it, are discussed in
\cite{devienne2011euclidean,layer2017embeddings}. The geometry of where the
principal axes cut a tree is treated in \cite{fiedler2005euclidean}.
\end{remark}

\begin{lemma}[Gap bracket]\label{d:lem:bracket}
$\displaystyle \max_{e\in E}w_e\;\le\;\lambda_1(\Mop)\;\le\;\sum_{e\in E}w_e$.
\end{lemma}

\begin{proof}
Lower: $\lambda_1(\Mop)\ge\ehat_e^\top\Mop\,\ehat_e\ge w_e$ for each $e$, since
every summand in \eqref{d:eq:splitcov} is positive semidefinite. Upper:
$\lambda_1(\Mop)\le\operatorname{tr}\Mop$, again by positive semidefiniteness.
\end{proof}

The inner products between split directions are also closed-form. Any two
distinct edges of a tree induce nested splits, so after orienting $A_f$
appropriately exactly one of the two cases below applies.

\begin{lemma}[Split Gram matrix]\label{d:lem:gram}
Let $e,f\in E$, $e\ne f$. Then
\begin{equation}\label{d:eq:gram}
  \langle\tilde u_e,\tilde u_f\rangle=
  \begin{cases}
    \dfrac{a_f\,b_e}{m}, & A_f\subseteq A_e,\\[2mm]
    -\dfrac{a_e\,a_f}{m}, & A_f\subseteq B_e .
  \end{cases}
\end{equation}
\end{lemma}

\begin{proof}
$\langle H\1_{A_e},H\1_{A_f}\rangle=\1_{A_e}^\top H\1_{A_f}
 =|A_e\cap A_f|-a_ea_f/m$. In the first case $|A_e\cap A_f|=a_f$, giving
$a_f(1-a_e/m)=a_fb_e/m$; in the second $|A_e\cap A_f|=0$.
\end{proof}

\Cref{d:eq:gram} is never zero, which is the reason the sign rule does not
simply return the edge of largest weight $w_e$: the directions $\ehat_e$ are not
orthogonal, and the top eigenvector of $\Mop$ responds to the whole correlated
family. \Cref{d:sec:empirical} records how often the two differ.

%==============================================================================
\section{Validity, unconditionally}\label{d:sec:validity}
%==============================================================================

The second structure theorem says that $-\tfrac12\Bop$ is not merely some
positive semidefinite matrix but a Laplacian pseudoinverse. It is asserted
without proof in \cite[\S4.4, \S4.8.1]{griffing2012connections}; we prove it,
because everything in this section rests on it.

Give each edge the conductance $c_e:=1/\ell_e$ and let $L\in\R^{V\times V}$,
$V=X\sqcup R$, be the weighted Laplacian of $T$. Order $V$ with the leaves
first,
\begin{equation}\label{d:eq:block}
  L=\begin{pmatrix} L_{XX} & L_{XR}\\ L_{RX} & L_{RR}\end{pmatrix},
  \qquad
  L^{*}:=L/L_{RR}=L_{XX}-L_{XR}L_{RR}^{-1}L_{RX}.
\end{equation}
$L_{RR}$ is a principal submatrix of a connected graph Laplacian obtained by
deleting a non-empty vertex set, hence positive definite, so $L^{*}$ is
well defined.

\begin{theorem}[Gower matrix as a Laplacian pseudoinverse]\label{d:thm:pinv}
\begin{equation}\label{d:eq:pinv}
  G:=-\tfrac12\,H\Dmat H=\bigl(L^{*}\bigr)^{\dagger}.
\end{equation}
\end{theorem}

\begin{proof}
\emph{Step 1: $L^{*}$ is the boundary current-to-potential map.} Let
$x\in\R^{X}$ with $\1^\top x=0$, and let $v=(v_X,v_R)$ solve $Lv=(x,0)$. The
second block gives $v_R=-L_{RR}^{-1}L_{RX}v_X$, and substituting into the first
gives $L^{*}v_X=x$. So injecting the current $x$ at the leaves and none at the
internal vertices produces leaf potentials $v_X$ related to $x$ by $L^{*}$.
Consequently the effective resistance between $i,j\in X$ in the network $(T,c)$
equals $(e_i-e_j)^\top (L^{*})^{\dagger}(e_i-e_j)$.

\emph{Step 2: effective resistance is the tree metric.} In a tree the $i$--$j$
path is unique, so resistors in series give effective resistance
$\sum_{e\in\text{path}}1/c_e=\sum_{e\in\text{path}}\ell_e=\Dmat_{ij}$. Hence
$\Dmat$ is exactly the resistance matrix of $L^{*}$.

\emph{Step 3: double centring inverts the resistance formula.} For any
symmetric positive semidefinite $L^{*}$ with kernel $\1$, write
$g:=\operatorname{diag}\bigl((L^{*})^{\dagger}\bigr)$. Expanding Step 1,
$\Dmat_{ij}=g_i+g_j-2(L^{*})^{\dagger}_{ij}$, i.e.
$\Dmat=\1g^\top+g\1^\top-2(L^{*})^{\dagger}$. Apply $H(\cdot)H$: the first two
terms are annihilated because $H\1=0$, so
$H\Dmat H=-2\,H(L^{*})^{\dagger}H$. Finally $(L^{*})^{\dagger}\1=0$ and
$(L^{*})^{\dagger}$ is symmetric, so $H(L^{*})^{\dagger}H=(L^{*})^{\dagger}$,
which is \eqref{d:eq:pinv}.
\end{proof}

Pseudoinversion fixes eigenvectors and inverts the non-zero eigenvalues.
Both $G$ and $L^{*}$ have kernel exactly $\operatorname{span}(\1)$
(\Cref{d:cor:nsd} for $G$; connectivity for $L^{*}$), so:

\begin{corollary}[Spectral correspondence]\label{d:cor:corr}
$G$ and $L^{*}$ have the same eigenvectors. The eigenvector of $G$ at its
largest eigenvalue --- equivalently, by \Cref{d:rem:sign}, the vector $\uone$ of
\Cref{d:sec:setup} --- is the Fiedler vector of $L^{*}$, and
$\lambda_{\max}(G)=1/\lambda_2(L^{*})$.
\end{corollary}

$L^{*}$ is the Laplacian of a weighted graph on the leaves obtained from $T$ by
eliminating internal vertices, and the sign pattern of its Fiedler vector is
governed by a theorem of Stone and Griffing
\cite[Thm.~3.3]{stone2009fiedler}: it partitions $X$ into $V_+$ and $V_-$ such
that, together with a suitable partition of the internal vertices, both parts
induce connected subgraphs of $T$. This is the Schur-complemented counterpart of
Fiedler's classical result that each nodal domain of a Fiedler vector induces a
connected subgraph \cite{fiedler1975property}; the behaviour of the later
eigenvectors under the same reduction is described in
\cite{griffing2013interlacing}. A bipartition of the leaves whose two sides
induce connected subtrees is a split of $T$. Hence:

\begin{corollary}[Validity of the sign rule]\label{d:cor:validity}
$\operatorname{sign}(\uone)$ is the bipartition $(A_{e},B_{e})$ induced by some
single edge $e\in E$, and $\uone$ has no zero coordinate. No quantitative
hypothesis on $T$ is required.
\end{corollary}

\begin{definition}[Dominant edge]\label{d:def:dominant}
By \Cref{d:cor:validity} there is an edge $\est\in E$ with
$\operatorname{sign}(\uone)=(A_{\est},B_{\est})$. Write
$\Ast:=A_{\est}$, $\Bst:=B_{\est}$, $a:=|\Ast|$, $b:=|\Bst|$,
$\ell_*:=\ell_{\est}$, $\wst:=w_{\est}=2\ell_*ab/m$, and
$\delta_*:=\min_i|\uone_i|>0$.
\end{definition}

This is the target of everything that follows. The task of the remaining
sections is not to identify $\est$ --- \Cref{d:cor:validity} has already
guaranteed that the exact rule returns a split of $T$ --- but to say how much of
$\Dmat$ may be discarded before the sampled rule stops returning
$(\Ast,\Bst)$.

\begin{remark}[Relation to the Fiedler vector of a similarity Laplacian]
\label{d:rem:notthesame}
\Cref{d:cor:validity} certifies that the distance operator returns \emph{a}
split. It does not assert that this split is the one a Laplacian built from a
similarity matrix returns, and in general it is not: the two operators are
different functions of the tree and they select different edges. The
quantitative reason is isolated in \Cref{d:prop:balance}.
\end{remark}

%==============================================================================
\section{Population landmarks}\label{d:sec:landmarks}
%==============================================================================

Three population quantities enter the sampling analysis. Each is a property of
the fully observed operator alone.

\begin{assumption}[Spectral gap]\label{d:asm:gap}
$\lambda_1(\Mop)$ is simple, and
$\Delta:=\lambda_1(\Mop)-\lambda_2(\Mop)\ \ge\ \gamma\,\wst$
for a constant $\gamma>0$.
\end{assumption}

\begin{assumption}[Delocalisation]\label{d:asm:deloc}
$\delta_*=\min_i|\uone_i|\ \ge\ c_\delta/\sqrt{m}$ for a constant $c_\delta>0$.
\end{assumption}

\begin{assumption}[Row separation]\label{d:asm:rowsep}
For every $i\in X$, with the convention that the sums omit $j=i$,
\begin{equation}\label{d:eq:rowsep}
  s_i:=\Bigl|\;
    \frac{1}{|\Bst|}\sum_{j\in \Bst}\Dmat_{ij}
    -\frac{1}{|\Ast|-1}\sum_{j\in \Ast}\Dmat_{ij}\;\Bigr|
  \;\ge\;s_0>0
  \quad(i\in\Ast),
\end{equation}
and symmetrically for $i\in\Bst$. Write $\kappa_{\mathrm{ref}}:=\dmax/s_0$.
\end{assumption}

\Cref{d:asm:deloc} is a normalisation, not a restriction on the tree: by
\eqref{d:eq:deloc} the split direction $\ehat_{\est}$ itself has
$\min_i|\ehat_i|=\sqrt{\min(a,b)/(m\max(a,b))}$, which is
$c_\delta/\sqrt m$ with $c_\delta=(1+\eta)^{-1/2}$ at imbalance
$\eta=\max(a,b)/\min(a,b)$. \Cref{d:asm:rowsep} says the same split is visible
in a single row of $\Dmat$; by \Cref{d:lem:path} the edge $\est$ alone
contributes exactly $\ell_*$ to $s_i$, so $s_0$ is the part of that
contribution that survives the other edges.

\paragraph{A family in which all three are exact.}
Let $T$ be the balanced binary tree on $m=2^k$ leaves in which the two edges at
the root have length $L$ and every other edge has length $g$, so that
$\Dmat_{ij}=2L+2g(k-1)$ across the root split and $2gd$ within, $d$ being the
depth of the meet. Then $\Ast,\Bst$ are the two halves,
$\wst=2\cdot(2L)\cdot\tfrac{m}{2}\cdot\tfrac{m}{2}/m=Lm$, and $\uone$ equals
$\ehat_{\est}$ exactly, so $\delta_*=1/\sqrt m$ and $c_\delta=1$. Numerically,
for $g=1$:

\begin{center}
\begin{tabular}{rrrrrrr}
\toprule
$m$ & $L$ & $|\lambda_1(\Mop)|$ & $|\lambda_2(\Mop)|$ & $\Delta$ & $\wst=Lm$
    & $\sqrt m\,\delta_*$\\
\midrule
 64 & 0.5 &   94.00 &  62.00 &   32.00 &   32.00 & 1.000\\
 64 & 1.0 &  126.00 &  62.00 &   64.00 &   64.00 & 1.000\\
 64 & 2.0 &  190.00 &  62.00 &  128.00 &  128.00 & 1.000\\
 64 & 4.0 &  318.00 &  62.00 &  256.00 &  256.00 & 1.000\\
256 & 0.5 &  382.00 & 254.00 &  128.00 &  128.00 & 1.000\\
256 & 1.0 &  510.00 & 254.00 &  256.00 &  256.00 & 1.000\\
256 & 2.0 &  766.00 & 254.00 &  512.00 &  512.00 & 1.000\\
256 & 4.0 & 1278.00 & 254.00 & 1024.00 & 1024.00 & 1.000\\
\bottomrule
\end{tabular}
\end{center}

\noindent
so $\Delta=\wst$ identically and \Cref{d:asm:gap} holds with $\gamma=1$, while
$\lambda_2(\Mop)=m-2$ is set by the internal edges alone. Here
$\dmax=2L+2g(k-1)$, and the two ratios that govern the rate are
$\kappa:=\dmax/\ell_*$ and $\kappa_{\mathrm{ref}}=\dmax/s_0$.

%==============================================================================
\section{Observation model and spectral noise}\label{d:sec:noise}
%==============================================================================

Each unordered pair is observed independently with probability $p$: let
$\Omega_{ij}=\Omega_{ji}\sim\mathrm{Bernoulli}(p)$ for $i<j$, independent across
pairs, $\Omega_{ii}=0$, and form the inverse-probability estimator
\begin{equation}\label{d:eq:ipw}
  \widehat{\Dmat}:=p^{-1}\,\Dmat\circ\Omega,
  \qquad
  \widehat{\Bop}:=H\,\widehat{\Dmat}\,H .
\end{equation}

\begin{lemma}[Unbiasedness]\label{d:lem:unbiased}
$\E\widehat{\Dmat}=\Dmat$ and $\E\widehat{\Bop}=\Bop$. Moreover, with
$Z:=\widehat{\Dmat}-\Dmat$,
\begin{equation}\label{d:eq:proj}
  \opnorm{\widehat{\Bop}-\Bop}=\opnorm{HZH}\le\opnorm{Z}.
\end{equation}
\end{lemma}

\begin{proof}
$\E\Omega_{ij}=p$ gives the first claim entrywise; $H$ is deterministic, so it
passes through the expectation and centring introduces no bias. For
\eqref{d:eq:proj}, $H$ is an orthogonal projector, so $\opnorm{H}=1$.
\end{proof}

The entries of $Z$ on and above the diagonal are independent and mean zero,
with
\begin{equation}\label{d:eq:zbounds}
  |Z_{ij}|\le \dmax/p,
  \qquad
  \operatorname{Var}(Z_{ij})=\Dmat_{ij}^2\,\frac{1-p}{p}\le\frac{\dmax^2}{p}.
\end{equation}

\begin{lemma}[Spectral error]\label{d:lem:opnorm}
There are absolute constants $C_0,C_1,C_2$ such that if $pm\ge C_0$ then
\begin{equation}\label{d:eq:opnorm}
  \opnorm{Z}\;\le\;2.01\,\dmax\sqrt{\frac{m}{p}}
  \qquad\text{with probability at least } 1-C_1e^{-C_2pm}.
\end{equation}
\end{lemma}

\begin{proof}
$Z$ is symmetric with independent, mean-zero entries on and above the diagonal,
bounded in absolute value by $\dmax/p$ and with variance at most $\dmax^2/p$ by
\eqref{d:eq:zbounds}. This is the hypothesis of
\cite[Thm.~8.4]{chatterjee2015usvt}, applied to the matrix $pZ/\dmax$, whose
entries are bounded by $1$ and have variance at most $p$.
\end{proof}

\begin{remark}[Which concentration inequality, and why]\label{d:rem:whichbound}
\Cref{d:lem:opnorm} uses \cite{chatterjee2015usvt} because it is stated for
exactly this ensemble: a symmetric matrix with independent entries on and above
the diagonal. A fully self-contained derivation is available from the matrix
Bernstein inequality \cite[Thm.~1.4]{tropp2012userfriendly} --- writing $Z$ as a
sum of $\binom m2$ independent symmetric rank-two terms, with scale $\dmax/p$
and variance proxy $m\dmax^2/p$, gives
$\opnorm{Z}\lesssim\dmax\bigl(\sqrt{m\log m/p}+\log m/p\bigr)$, which costs a
factor $\sqrt{\log m}$. Two neighbouring results do not apply as stated:
\cite[Thm.~2]{achlioptas2007fast} is the algorithmic precedent for the same
$\sqrt{m/p}$ rate but is stated for the non-symmetric independent-entry model,
and \cite[Thm.~1.1]{bandeira2016sharp} is sharp but Gaussian.
\end{remark}

%==============================================================================
\section{From the spectrum to the labels}\label{d:sec:perturb}
%==============================================================================

Let $\huone$ be a unit eigenvector of $\widehat{\Bop}$ at its
largest-magnitude eigenvalue, sign-aligned so that
$\langle\huone,\uone\rangle\ge0$.

\begin{lemma}[$\ell_2$ control]\label{d:lem:dk}
Under \Cref{d:asm:gap}, on the event of \Cref{d:lem:opnorm},
\begin{equation}\label{d:eq:dk}
  \norm{\huone-\uone}_2\ \le\ \sqrt2\,\sin\Theta(\huone,\uone)
  \ \le\ \frac{2\sqrt2\,\opnorm{Z}}{\Delta}
  \ \le\ \frac{5.69\,\dmax}{\gamma\,\wst}\sqrt{\frac{m}{p}} .
\end{equation}
In particular, if $\Ast$ and $\Bst$ are balanced up to a constant, so that
$\wst\asymp\ell_*m$, then with $\kappa=\dmax/\ell_*$,
\begin{equation}\label{d:eq:dk2}
  \norm{\huone-\uone}_2\ \lesssim\ \frac{\kappa}{\gamma\sqrt{pm}} .
\end{equation}
\end{lemma}

\begin{proof}
Apply the one-dimensional case of the statistical variant
\cite[Cor.~1]{yu2015useful} of the Davis--Kahan $\sin\Theta$ theorem
\cite{davis1970rotation} to the symmetric pair $\Mop$, $-\widehat{\Bop}$,
whose difference has operator norm at most $\opnorm{Z}$ by
\eqref{d:eq:proj}; the corollary requires only the \emph{population} gap
$\Delta$, which is \Cref{d:asm:gap}. The bound
$\norm{\huone-\uone}_2\le\sqrt2\sin\Theta$ holds after sign alignment. Substitute
\eqref{d:eq:opnorm} and $\Delta\ge\gamma\wst$.
\end{proof}

\begin{proposition}[Misclassification count]\label{d:prop:misclass}
Let $\mathcal{E}:=\{i:\operatorname{sign}(\huone_i)\ne\operatorname{sign}(\uone_i)\}$.
Under \Cref{d:asm:gap,d:asm:deloc},
\begin{equation}\label{d:eq:misclass}
  \frac{|\mathcal{E}|}{m}\ \le\ \frac{\norm{\huone-\uone}_2^2}{m\,\delta_*^2}
  \ \le\ \frac{1}{c_\delta^2}\,\norm{\huone-\uone}_2^2
  \ \lesssim\ \frac{\kappa^2}{c_\delta^2\gamma^2\,pm}.
\end{equation}
\end{proposition}

\begin{proof}
If $i\in\mathcal{E}$ then $\huone_i$ and $\uone_i$ have opposite signs, so
$|\huone_i-\uone_i|\ge|\uone_i|\ge\delta_*$. Hence
$|\mathcal{E}|\,\delta_*^2\le\norm{\huone-\uone}_2^2$. Insert
\Cref{d:asm:deloc} and \eqref{d:eq:dk2}.
\end{proof}

\begin{remark}[Why not carry the eigenvector to $\ell_\infty$ directly]
\label{d:rem:linf}
Exact recovery by the raw sign rule needs
$\norm{\huone-\uone}_\infty<\delta_*$. The only dimension-free transfer,
$\norm{\cdot}_\infty\le\norm{\cdot}_2$, charges every leaf as though it were the
worst one, and against $\delta_*\asymp m^{-1/2}$ that costs a factor $m$: it
turns \eqref{d:eq:dk2} into the requirement $p\gtrsim\kappa^2/(c_\delta\gamma)^2$,
a constant sampling rate. \Cref{d:sec:exact} recovers the $\log m/m$ rate by
adding a second stage instead; \Cref{d:rem:entrywise} names the alternative.
\end{remark}

%==============================================================================
\section{Exact recovery by one-step refinement}\label{d:sec:exact}
%==============================================================================

Let $(\widehat{C}_0,\widehat{C}_1)$ be the partition
$\operatorname{sign}(\huone)$ from \Cref{d:sec:perturb}, and define the refined
labelling: assign $i$ to the part $c$ minimising
\begin{equation}\label{d:eq:refine}
  \psi_c(i):=\frac{1}{|\widehat{C}_c\setminus\{i\}|}
             \sum_{j\in \widehat{C}_c\setminus\{i\}}\widehat{\Dmat}_{ij},
  \qquad c\in\{0,1\}.
\end{equation}

\begin{lemma}[Per-leaf concentration]\label{d:lem:pernode}
Fix $i$ and a set $C\subseteq X\setminus\{i\}$ with $|C|\ge m/4$. For any
$\tau\in(0,1)$,
\begin{equation}\label{d:eq:bernstein}
  \Prob\Bigl(\bigl|\psi_C(i)-\bar\Dmat_{iC}\bigr|
       \ge \dmax\sqrt{\frac{8\log(m/\tau)}{p\,|C|}}
           +\frac{4\dmax\log(m/\tau)}{3p\,|C|}\Bigr)\le\frac{2\tau}{m},
\end{equation}
where $\bar\Dmat_{iC}:=|C|^{-1}\sum_{j\in C}\Dmat_{ij}$.
\end{lemma}

\begin{proof}
$\psi_C(i)=|C|^{-1}\sum_{j\in C}p^{-1}\Omega_{ij}\Dmat_{ij}$ is an average of
$|C|$ independent random variables with mean $\bar\Dmat_{iC}$, each bounded in
absolute value by $\dmax/(p|C|)$ after scaling and with variance at most
$\dmax^2/(p|C|^2)$ by \eqref{d:eq:zbounds}. Bernstein's inequality for bounded
independent summands gives \eqref{d:eq:bernstein}.
\end{proof}

\begin{theorem}[Exact recovery]\label{d:thm:main}
Assume \Cref{d:asm:gap,d:asm:deloc,d:asm:rowsep}, and that the imbalance
$\eta=\max(a,b)/\min(a,b)$ is bounded. There are constants
$C=C(\gamma,c_\delta,\eta)$ and $c>0$ such that if
\begin{equation}\label{d:eq:mainrate}
  p\;\ge\;C\,\bigl(\kappa_{\mathrm{ref}}^2+\kappa^2\bigr)\,
              \frac{\log m}{m},
  \qquad \kappa=\frac{\dmax}{\ell_*},\quad
         \kappa_{\mathrm{ref}}=\frac{\dmax}{s_0},
\end{equation}
then with probability at least $1-m^{-c}$ the refined labelling equals
$\operatorname{sign}(\uone)$ at every leaf. By \Cref{d:cor:validity} that
labelling is the bipartition of $X$ induced by the edge $\est$ of $T$.
\end{theorem}

\begin{proof}
Write $\varepsilon:=|\mathcal{E}|/m$ for the stage-one error rate. By
\Cref{d:prop:misclass} and \eqref{d:eq:mainrate},
$\varepsilon\lesssim\kappa^2/(c_\delta\gamma)^2pm\le s_0/(8\dmax)$ once $C$ is
large enough; in particular $\varepsilon<\tfrac12\min(a,b)/m$, so both
$\widehat{C}_0$ and $\widehat{C}_1$ have at least $m/4$ elements and each is
within $\varepsilon m$ leaves of the corresponding true part.

Fix $i\in\Ast$. Replacing $\widehat{C}_0,\widehat{C}_1$ by $\Ast,\Bst$ inside
\eqref{d:eq:refine} changes each average by at most
$2\varepsilon\dmax\cdot(m/\min|\widehat{C}_c|)\le 8\varepsilon\dmax\le s_0/4$,
because at most $\varepsilon m$ summands differ and each is bounded by $\dmax$.
By \Cref{d:asm:rowsep} the population difference
$\bar\Dmat_{i\Bst}-\bar\Dmat_{i\Ast}$ has absolute value at least $s_0$ and the
sign that identifies $i$ as belonging to $\Ast$. It therefore suffices that the
stochastic deviation of each of the two averages be below $s_0/4$. By
\Cref{d:lem:pernode} with $\tau=m^{-c}$ and $|C|\ge m/4$, both deviations are at
most
\[
  \dmax\sqrt{\frac{32(1+c)\log m}{p\,m}}
  +\frac{16(1+c)\dmax\log m}{3pm}
  \;\le\;\frac{s_0}{4}
\]
whenever $pm\ge C'\kappa_{\mathrm{ref}}^2\log m$, which is
\eqref{d:eq:mainrate}. A union bound over the $2m$ statistics
$\psi_c(i)$ costs a factor $2m$, absorbed by $\tau=m^{-c}$. The argument for
$i\in\Bst$ is identical.
\end{proof}

\begin{corollary}[Rate]\label{d:cor:rate}
If $\gamma$, $c_\delta$, $\eta$, $\kappa$ and $\kappa_{\mathrm{ref}}$ are
bounded as $m$ grows, then $p=\Theta(\log m/m)$ suffices, and the expected number
of observed entries is $\Theta(m\log m)$ --- a vanishing fraction
$\Theta(\log m/m)$ of the $\binom m2$ pairwise distances.
\end{corollary}

\begin{remark}[Claim discipline]\label{d:rem:sufficient}
\eqref{d:eq:mainrate} is a sufficient condition. No matching lower bound is
established here, so $\Theta$ in \Cref{d:cor:rate} names the growth order of a
sufficient rate, not the optimal one.
\end{remark}

\begin{remark}[The single-stage rule]\label{d:rem:singlestage}
Stage one alone --- the raw sign rule on $\huone$, with no refinement --- is
covered by \Cref{d:prop:misclass}, which gives an error rate
$O(\kappa^2/(c_\delta\gamma)^2pm)$, and by \Cref{d:rem:linf}, which gives exact
recovery at the constant rate $p\gtrsim\kappa^2/(c_\delta\gamma)^2$.
\end{remark}

\begin{remark}[Entrywise alternative]\label{d:rem:entrywise}
The factor $m$ lost in \Cref{d:rem:linf} is an artefact of routing an
$\ell_\infty$ event through an $\ell_2$ bound. A leave-one-out decomposition of
$\huone$ bounds $\norm{\huone-\uone}_\infty$ directly
\cite{abbe2020entrywise,chen2021spectral} and yields exact recovery from the
single-stage rule at the rate of \eqref{d:eq:mainrate}. The two-stage argument
above is preferred here because \Cref{d:lem:pernode} is elementary and its
hypotheses are visible. Two-stage refinement of a spectral initialiser is the
standard device for converting an $\ell_2$ guarantee into exact labelling
\cite{gao2017achieving,loffler2021optimality}.
\end{remark}

%==============================================================================
\section{What the operator selects}\label{d:sec:empirical}
%==============================================================================

\Cref{d:thm:splitcov} contains a structural statement about \emph{which} split
the operator favours, independent of any sampling.

\begin{proposition}[Balance preference]\label{d:prop:balance}
For a fixed edge length $\ell$, the weight $w_e=2\ell a_eb_e/m$ attached to a
split of sizes $(a_e,b_e)$ satisfies
\begin{equation}\label{d:eq:balance}
  w_e=\frac{\ell m}{2}\cdot\frac{4\eta_e}{(1+\eta_e)^2},
  \qquad \eta_e:=\frac{\max(a_e,b_e)}{\min(a_e,b_e)},
\end{equation}
which is maximal at $\eta_e=1$ and decays like $4/\eta_e$ as
$\eta_e\to\infty$. An edge must therefore be longer by a factor
$\asymp\eta_e/4$ to compete with a balanced edge of the same length.
\end{proposition}

\begin{proof}
Substitute $a_e=m/(1+\eta_e)$, $b_e=m\eta_e/(1+\eta_e)$ into $2\ell a_eb_e/m$.
\end{proof}

Two consequences are visible in simulation. First, the splits returned by the
sign rule are markedly more balanced than those returned by a Fiedler vector of
a similarity Laplacian, which carries no analogue of the factor $a_eb_e$.
Across $1000$ trees per cohort at $m=1000$, with no sub-sampling and each
operator using the rounding rule its own analysis states:

\begin{center}
\begin{tabular}{lrrrr}
\toprule
 & \multicolumn{2}{c}{Kingman} & \multicolumn{2}{c}{birth--death}\\
\cmidrule(lr){2-3}\cmidrule(lr){4-5}
 & $L(S)$ & $\Bop$ & $L(S)$ & $\Bop$\\
\midrule
valid single-edge split & 996 & 986 & 996 & 555\\
median $\eta$           & 3.00 & 1.85 & 3.08 & 1.53\\
maximum $\eta$          & 499 & 24.6 & 999 & 5.3\\
\bottomrule
\end{tabular}
\end{center}

\noindent
The median and maximum imbalance of $\Bop$ are bounded where those of $L(S)$ are
not, which is \eqref{d:eq:balance}. The same pattern appears on random weighted
trees generated without a coalescent: over $200$ trees at $m=64$ with
independent uniform edge lengths, $\Bop$ returned a valid single-edge split in
$200$ cases out of $200$ --- \Cref{d:cor:validity} --- with median $\eta=1.29$
and maximum $\eta=4.8$.

Second, on those same $200$ trees the selected split was the one of largest
weight $w_e$ in only $40\%$ of cases, the median rank being $2$. This is
\Cref{d:lem:gram} at work: the directions $\ehat_e$ are mutually correlated, so
the top eigenvector of $\Mop$ is not the single heaviest projector but the
direction that best explains the whole family.

The birth--death column reads the same way through \Cref{d:asm:gap}. Under that
cohort the products $\ell_ea_eb_e$ are more nearly equal across edges, $\Delta$
is small relative to $\wst$, and $\uone$ is a mixture of several $\ehat_e$ whose
sign pattern need not be any single split --- the regime in which
\Cref{d:asm:gap} is the operative hypothesis and in which
\eqref{d:eq:mainrate} degrades through $\gamma$.

%==============================================================================
\section{Scope, and the two routes side by side}\label{d:sec:scope}
%==============================================================================

The results above are stated for additive metrics of a weighted tree, for the
pairwise Bernoulli observation model \eqref{d:eq:ipw}, and under
\Cref{d:asm:gap,d:asm:deloc,d:asm:rowsep} with bounded
$\kappa,\kappa_{\mathrm{ref}},\gamma,c_\delta,\eta$. Within that scope the
conclusion is exact recovery of the fully observed operator's own partition, at
$p=\Theta(\log m/m)$, from $\Theta(m\log m)$ observed distances.

\Cref{d:cor:validity} is stronger in one respect and weaker in another than the
quantitative results: it holds for every weighted tree with no hypothesis at
all, and it certifies that the partition is a split of $T$ --- but it says which
split only through \Cref{d:def:dominant}, i.e.\ by naming the edge the exact
operator selects. That is the appropriate form of the claim, because the
distance operator and a similarity Laplacian select different edges by
construction (\Cref{d:prop:balance}), and a guarantee that pinned the distance
route to the other operator's edge would be false.

The four steps correspond to those of the similarity route as follows.

\begin{center}
\begin{tabular}{llll}
\toprule
Step & Similarity route & Distance route & Where they differ\\
\midrule
1. Population geometry
  & block structure of $L$
  & \Cref{d:thm:splitcov,d:thm:pinv}
  & validity is unconditional here\\
2. Operator-norm error
  & matrix Bernstein
  & \Cref{d:lem:opnorm}
  & symmetric-ensemble bound, no $\sqrt{\log m}$\\
3. Eigenvector control
  & resolvent expansion
  & \Cref{d:lem:dk}
  & Davis--Kahan; no series to truncate\\
4. Per-node concentration
  & scalar Bernstein
  & \Cref{d:lem:pernode}
  & same tool, applied after refinement\\
\bottomrule
\end{tabular}
\end{center}

\noindent
Step 3 is the deliberate divergence. A resolvent expansion produces an
$\ell_\infty$ bound in one stroke and so needs no second stage, at the price of
controlling the tail of the series, which requires the ratio
$\opnorm{Z}/\Delta$ to be small. \Cref{d:lem:dk} asks nothing of that ratio and
pushes the $\ell_\infty$ work into \Cref{d:lem:pernode}, where it is a scalar
Bernstein bound over the $\asymp pm$ entries actually observed in one row.

\begin{remark}[Reproducibility of the numerical statements]\label{d:rem:numerics}
Every identity asserted without a displayed proof constant ---
\eqref{d:eq:path}, \eqref{d:eq:centring}, \eqref{d:eq:splitcov},
\eqref{d:eq:trace}, \eqref{d:eq:gram} and \eqref{d:eq:pinv} --- was checked
numerically on random weighted trees to a maximum absolute error of
$1.4\times10^{-14}$, as was the landmark table of \Cref{d:sec:landmarks} and the
scaling $p\propto\kappa^2\log m/m$ of \Cref{d:thm:main} over
$m=256,\dots,2048$.
\end{remark}

\bibliography{refs}

\end{document}
